% !TeX root = ../../Book1.tex
\section{极大不等式的基本应用}
\label{b1:sec:1304}

极大不等式的用途不止是估计一个新函数的范数。它把“一族平均中至少有一个偏大”这样的条件，变成统一的例外集估计；这正是从范数逼近走向几乎处处结论时缺少的信息。本节先明确允许哪些集合趋近一个点，再讨论误差控制与积分核。下一章才证明局部平均恢复函数值的微分定理，这里把其中会反复使用的工具准备好。

\subsection{微分基}
\label{b1:subsec:130401}

\begin{definition}\label{b1:def:differentiation-basis}
对每个 \(x\in\mathbb R^d\)，给定一族包含 \(x\) 的有界 Lebesgue 可测集 \(\mathcal B(x)\)。若其中每个 \(E\) 都满足 \(0<m_d(E)<\infty\)，且对每个 \(\delta>0\) 都存在 \(E\in\mathcal B(x)\) 使 \(\operatorname{diam}E<\delta\)，则称 \(\mathcal B\) 为一个\term[weifen-ji]{微分基}[differentiation basis]。
若对给定的局部可积函数 \(f\)，有
\[
 \lim_{\substack{E\in\mathcal B(x)\\\operatorname{diam}E\to0}}
       \frac1{m_d(E)}\int_E f(y)\,\mathrm dy=f(x)
\]
几乎处处成立，就说此基能微分 \(f\)。上式的极限要求所有足够小的允许集合都满足同一个误差界，不只是选定某个集合序列。
\end{definition}

这里的“基”只描述允许集合怎样趋近一点；它是否能够微分某个函数类，是需要另外证明的性质，不能从定义中的任意小尺度条件直接推出。

\ProseParagraphBreak

中心球、含点球和以 \(x\) 为中心的立方体都给出微分基。对连续函数，它们都能恢复函数值：只要 \(x\in E\) 且 \(\operatorname{diam}E<\delta\)，就有
\[
 \left|\frac1{m_d(E)}\int_E f-f(x)\right|
 \leq\sup_{|y-x|<\delta}|f(y)-f(x)|\longrightarrow0.
\]
对不连续的可积函数，连续性不能提供这个上界，允许集合的形状便开始影响结论。

\ProseParagraphBreak

\cref{b1:fig:differentiation-bases}把三种常见的趋近方式并列画出。中心条件控制球心与观察点的关系；任意长宽比的矩形则可能在一个方向很短、另一个方向仍相对很长。图中只比较几何限制，不预先断言哪一族能微分全部局部可积函数。

\begin{figure}[htbp]
\centering
\begin{tikzpicture}[scale=0.95,line cap=round,line join=round]
  \begin{scope}[xshift=-4.3cm]
    \draw[thick] (0,0) circle (1);
    \fill (0,0) circle (1.5pt);
    \node[above right] at (0,0) {$x$};
    \node[below] at (0,-1.35) {中心球};
  \end{scope}
  \begin{scope}
    \draw[thick] (0.35,0.05) circle (1);
    \draw (0.35,0.05) circle (1.5pt);
    \fill (-0.25,0.30) circle (1.5pt);
    \node[above left] at (-0.25,0.30) {$x$};
    \node[below right] at (0.35,0.05) {$c$};
    \node[below] at (0,-1.35) {含点球};
  \end{scope}
  \begin{scope}[xshift=4.3cm]
    \draw[thick] (-1.65,-0.28) rectangle (1.65,0.28);
    \fill (0,0) circle (1.5pt);
    \node[above] at (0,0) {$x$};
    \node[below] at (0,-1.35) {任意长宽比的矩形};
  \end{scope}
\end{tikzpicture}
\caption[三类微分基]{三类微分基。前两类保持球形，但观察点与球心的关系不同；第三类若不限制长宽比，偏心率可以任意增大。}
\label{b1:fig:differentiation-bases}
\end{figure}

一种可用的几何条件是：对每个 \(E\in\mathcal B(x)\)，都能找到 \(r>0\)，使
\begin{equation}\label{b1:eq:basis-volume-comparison}
 E\subseteq B(x,r),\quad m_d\bigl(B(x,r)\bigr)\leq A\,m_d(E),
\end{equation}
其中 \(A\) 不依赖 \(x,E\)。于是对每个局部可积 \(f\)，
\[
 \frac1{m_d(E)}\int_E|f|\leq A Mf(x).
\]
如果同时保证所选 \(r\) 随 \(\operatorname{diam}E\) 趋零，就既控制了平均的大小，也保留了趋向该点的尺度。中心立方体满足这种条件；边长比例任意失衡的长方体却不能取统一的 \(A\)。本章习题将给出后一类形状的极大弱型估计失败的具体例子。

\ProseParagraphBreak

\secref{b1:sec:1301}的覆盖关系记录点与集合的配对，本节的微分基则额外要求平均的分母为有限正数，并把趋零极限作为研究对象。仅仅满足微分基的定义，不等于已经具有 Vitali 选择性质，也不等于能够微分全部局部可积函数。

\subsection{\texorpdfstring{\(L^p\)}{Lp} 控制}
\label{b1:subsec:130402}

若 \(A_rf(x)\) 表示中心球上不取绝对值的平均，则
\[
 |A_rf(x)-A_rg(x)|\leq M(f-g)(x)
\]
对所有 \(r>0\) 同时成立。因此，\(f-g\) 的范数虽只给出整体误差，极大不等式却控制了“存在某个尺度使平均误差超过阈值”的位置。在 \(L^1\) 中用弱型估计，在 \(L^p\)、\(p>1\) 中还可用强型估计。

\ProseParagraphBreak

下面把这一方法抽成一个序列形式。它不预先断言某种具体微分基已经有效。

% 思路来源：实读Axler2020Measure，4.10的连续稠密类加极大误差方法；
% 改编为可数线性算子族上的闭包原则，证明及常数单独推导。
\begin{proposition}[极大估计下的收敛延拓]\label{b1:prop:maximal-convergence-principle}
设 \(1\leq p<\infty\)，\(T_n\) 为从 \(L^p(\mu)\) 到有限值可测函数的线性算子，按几乎处处相等理解。令
\[
 T^*f=\sup_{n\geq1}|T_nf|.
\]
假设存在 \(A<\infty\)，对每个 \(f\in L^p\)、\(t>0\)，有
\[
 \mu\bigl(\{T^*f>t\}\bigr)\leq\frac A{t^p}\|f\|_p^p.
\]
若有一个在 \(L^p\) 中稠密的子集 \(\mathcal D\)，使 \(T_ng\to g\) 对每个 \(g\in\mathcal D\) 几乎处处成立，则 \(T_nf\to f\) 对所有 \(f\in L^p\) 几乎处处成立。
\end{proposition}
\begin{proof}
固定 \(f\) 与 \(g\in\mathcal D\)。选定可测代表，删去与它们有关的可数个线性等式失效点以及 \(T_ng\not\to g\) 的零测集。在其余点，
\[
 \limsup_{n\to\infty}|T_nf-f|
 \leq T^*(f-g)+|f-g|.
\]
因此对 \(t>0\)，极大估计和 Chebyshev 不等式给出
\[
 \begin{aligned}
 \mu\left(\left\{\limsup_n|T_nf-f|>2t\right\}\right)
 &\leq\mu\bigl(\{T^*(f-g)>t\}\bigr)
       +\mu\bigl(\{|f-g|>t\}\bigr)\\
 &\leq\frac{A+1}{t^p}\|f-g\|_p^p.
 \end{aligned}
\]
左端不依赖 \(g\)。由稠密性让右端任意小，可知左端为零。再取 \(t=1/k\)、\(k\geq1\)，作可数并，即得几乎处处收敛。
\end{proof}

证明并没有让逼近序号与算子序号同时趋于无穷，而是先对一个固定的稠密类元素取算子极限，再把剩余误差送到零。极大函数负责统一所有算子序号，因此不需要对每个 \(n\) 另选一组例外点。该命题只处理可数族；连续半径情形还应说明可测性或可数尺度的约化，不能把不可数个零测例外集直接并起来。

\subsection{卷积与 \texorpdfstring{\(L^p\)}{Lp} 平移连续性}
\label{b1:subsec:1304-convolution-translation}

球平均可以写成一个固定函数的缩放与积分。为了容纳更平滑的平均方式，本小节建立卷积的 \(L^p\) 有界性和平移连续性；光滑性与高阶导数留到第四编处理。

\ProseParagraphBreak

给定可测函数 \(K,f\)，在绝对积分有限的点定义它们的\term[juanji]{卷积}[convolution]为
\[
 (K*f)(x)\coloneqq\int_{\mathbb R^d}K(y)\cdot f(x-y)\,\mathrm dy.
\]
\symidx{\(K*f\)：函数卷积}
取 Borel 代表后被积函数关于 \((x,y)\) 可测，换回几乎处处相等的代表不改变积分。事实上，固定 \(x\) 时，反射和平移都保持 Lebesgue 零测集。下文把积分不存在的零测集上的值约定为零；涉及某个具体点的结论时，则另行证明该点的积分存在。卷积的这些测度论基础可参看 Fremlin 的《Measure Theory》\cite[255F--L]{Fremlin2016Measure}。

\begin{lemma}\label{b1:lem:convolution-lp-bound}
若 \(K\in L^1(\mathbb R^d)\)、\(f\in L^p(\mathbb R^d)\)、\(1\leq p\leq\infty\)，则 \(K*f\) 几乎处处绝对存在，且
\[
 \|K*f\|_p\leq\|K\|_1\cdot\|f\|_p.
\]
\end{lemma}
\begin{proof}
令 \(a=\|K\|_1\)。若 \(a=0\)，结论直接成立。设 \(a>0\)、\(p<\infty\)。对概率测度 \(|K(y)|\,\mathrm dy/a\) 应用 Hölder 不等式，得
\[
 \left(\int|K(y)|\cdot|f(x-y)|\,\mathrm dy\right)^p
 \leq a^{p-1}\int|K(y)|\cdot|f(x-y)|^p\,\mathrm dy.
\]
当 \(p=1\) 时，这是等式。右侧对 \(x\) 积分，由 Tonelli 定理和平移不变性得 \(a^p\|f\|_p^p<\infty\)。所以绝对积分几乎处处有限；再用积分三角不等式即可得到范数界。若 \(p=\infty\)，选择处处有界代表后，绝对积分直接受 \(a\|f\|_\infty\) 控制。
\end{proof}

上一引理处理的是一个因子可积时的端点。后文还会需要两个有限指数共同决定输出指数的版本，因此在卷积理论首次建立处一并记录有限指数的\term[young-juanji-budengshi]{Young 卷积不等式}[Young convolution inequality]。它与\secref{b1:sec:1001}的标量 Young 不等式有关，但陈述对象不同。

\begin{lemma}[Young]\label{b1:lem:young-convolution-finite}
设 \(1\le a,b,q<\infty\)，满足
\[
 1+\frac1q=\frac1a+\frac1b.
\]
若 \(f\in L^a(\mathbb R^d)\)、\(g\in L^b(\mathbb R^d)\)，则卷积几乎处处绝对收敛，并有
\[
 \|f*g\|_{L^q}\leq\|f\|_{L^a}\cdot\|g\|_{L^b}.
\]
\end{lemma}
\begin{proof}
条件保证 \(q\ge a,b\)。对固定 \(x\)，把绝对值乘积写为三因子：
\[
 |f(x-y)|\cdot|g(y)|
 =\bigl(|f(x-y)|^a|g(y)|^b\bigr)^{1/q}
    |f(x-y)|^{1-a/q}|g(y)|^{1-b/q}.
\]
Hölder 不等式所用的三个倒指数为
\[
 \frac1q,\quad \frac1a-\frac1q,\quad \frac1b-\frac1q,
\]
其和为一；若某一后项为零，就删去对应的零幂因子。由此
\[
 \left(\int|f(x-y)|\cdot|g(y)|\,\mathrm dy\right)^q
 \leq \|f\|_a^{q-a}\cdot\|g\|_b^{q-b}
         \int |f(x-y)|^a|g(y)|^b\,\mathrm dy.
\]
对 \(x\) 积分并使用 Tonelli 定理及平移不变性，右侧等于
\(\|f\|_a^q\cdot\|g\|_b^q\)。因此左侧有限，这同时证明几乎处处绝对收敛和所求界。零范数情形直接为零，\(a=b=q=1\) 时整个论证就是 Tonelli 定理。
\end{proof}

\begin{lemma}\label{b1:lem:lp-translation-continuity}
若 \(1\leq p<\infty\)、\(f\in L^p(\mathbb R^d)\)，记 \(\tau_hf(x)=f(x-h)\)，则
\[
 \|\tau_hf-f\|_p\longrightarrow0\quad(h\to0).
\]
\end{lemma}
\begin{proof}
由\cref{b1:prop:lp-continuous-dense}，给定 \(\varepsilon>0\)，可取 \(g\in C_c(\mathbb R^d)\) 使 \(\|f-g\|_p<\varepsilon\)。当 \(|h|\leq1\) 时，\(\tau_hg-g\) 的支集包含在一个固定紧集 \(H\) 中。连续紧支集函数在整个空间一致连续，故
\[
 \|\tau_hg-g\|_p
 \leq m_d(H)^{1/p}\sup_x|g(x-h)-g(x)|\longrightarrow0.
\]
平移保持 \(L^p\) 范数，所以
\[
 \|\tau_hf-f\|_p
 \leq2\|f-g\|_p+\|\tau_hg-g\|_p.
\]
先取 \(h\to0\)，再让 \(\varepsilon\to0\)，得到结论。
\end{proof}

这一步使用的是\chapref{b1:ch:10}已经建立的连续紧支集稠密性，不预借光滑函数稠密性。当 \(p=\infty\) 时结论不成立：对 \(f=\mathbf1_{(0,1)}\) 和 \(0<|h|<1\)，有 \(\|\tau_hf-f\|_\infty=1\)。

\subsection{逼近恒等与极大函数控制}
\label{b1:subsec:130403}

有了卷积的范数估计和平移连续性，可以先证明固定核缩放在 \(L^p\) 中逼近恒等，再说明径向递减核为何同时受 Hardy--Littlewood 极大函数控制。

\begin{definition}\label{b1:def:approximate-identity}
设 \(K\in L^1(\mathbb R^d)\) 且 \(\int K=1\)。令
\[
 K_t(x)\coloneqq t^{-d}\cdot K(x/t),\quad t>0.
\]
称 \((K_t)_{t>0}\) 为由 \(K\) 生成的一族\term[bijin-hengdeng]{逼近恒等}[approximate identity]。这里采用固定核缩放的版本，不要求 \(K\) 非负或光滑。
\end{definition}
\symidx{\(K_t\)：积分核的归一化缩放}

对每个 \(t\)，\(\int K_t=1\)、\(\|K_t\|_1=\|K\|_1\)；而对每个 \(\delta>0\)，
\[
 \int_{|x|>\delta}|K_t(x)|\,\mathrm dx
 =\int_{|z|>\delta/t}|K(z)|\,\mathrm dz\longrightarrow0.
\]
所以核的积分保持不变，绝对积分的质量则集中到原点附近。这个描述指测度意义下的集中，不是 \(K_t\) 在原点附近一致有界。

\begin{proposition}\label{b1:prop:approximate-identity-lp}
对上述 \(K_t\) 及 \(f\in L^p(\mathbb R^d)\)、\(1\leq p<\infty\)，有
\[
 \|K_t*f-f\|_p\longrightarrow0\quad(t\downarrow0).
\]
若 \(f\) 有界且一致连续，则 \(K_t*f\to f\) 一致成立。
\end{proposition}
\begin{proof}
作代换 \(y=tz\)，并用 \(\int K=1\)，得到
\[
 K_t*f(x)-f(x)=\int K(z)\bigl(f(x-tz)-f(x)\bigr)\,\mathrm dz
\]
几乎处处成立。令 \(a=\|K\|_1\geq1\)，重复前一卷积界中的概率测度估计，再对 \(x\) 积分，可得
\[
 \|K_t*f-f\|_p^p
 \leq a^{p-1}\int |K(z)|\cdot\|\tau_{tz}f-f\|_p^p\,\mathrm dz.
\]
对每个固定 \(z\)，平移连续性使右侧的范数趋零；同时它不超过 \(2^p\|f\|_p^p\)。控制收敛定理因而完成有限 \(p\) 的证明。

若 \(f\) 有界且一致连续，定义其平移模
\[
 \omega_f(h)=\sup_x|f(x-h)-f(x)|.
\]
由平移变量和三角不等式，
\[
 \omega_f(-h)=\omega_f(h),
 \quad
 |\omega_f(h)-\omega_f(k)|\leq\omega_f(h-k).
\]
一致连续性给出 \(\omega_f(u)\to0\) 当 \(u\to0\)，所以上式又说明 \(\omega_f\) 在每一点连续，特别地是 Borel 可测函数。现在同一个卷积表示给出
\[
 \sup_x|K_t*f(x)-f(x)|
 \leq\int|K(z)|\omega_f(tz)\,\mathrm dz.
\]
对每个固定 \(z\)，\(\omega_f(tz)\to0\)，而且
\(0\leq\omega_f(tz)\leq2\|f\|_\infty\)。因此被积函数可测，并由可积函数 \(2\|f\|_\infty|K(z)|\) 控制；控制收敛定理给出一致收敛。
\end{proof}

上述范数收敛不需要极大函数，但它本身不能把整个连续参数族的极限变成每一点的极限。为获得逐点信息，还需一种能统一所有 \(t\) 的控制。径向递减核正好提供了这样的联系。

\begin{proposition}\label{b1:prop:radial-kernel-maximal}
设 \(K\geq0\)、\(K\in L^1(\mathbb R^d)\)，且有一个径向非增代表，即 \(K(x)=k(|x|)\)，其中 \(k\) 非增。对每个 \(f\in L^1_{\mathrm{loc}}\)、每个 \(t>0\)，有扩展实数意义的不等式
\[
 \int K_t(y)|f(x-y)|\,\mathrm dy\leq\|K\|_1Mf(x).
\]
因此在 \(Mf(x)<\infty\) 的每个点，所有 \(K_t*f(x)\) 均绝对存在，并且
\[
 \sup_{t>0}|K_t*f(x)|\leq\|K\|_1Mf(x).
\]
\end{proposition}
\begin{proof}
固定 \(t\)。对 \(s>0\)，集合 \(E_s=\{y:K_t(y)>s\}\) 除去一个球面后是以原点为中心的球，也可能为空或只含原点。它具有有限测度，因为 \(s\,m_d(E_s)\leq\|K\|_1\)。故平移及球平均的定义给出
\[
 \int_{E_s}|f(x-y)|\,\mathrm dy\leq m_d(E_s)\cdot Mf(x).
\]
若 \(Mf(x)=\infty\) 且 \(\|K\|_1>0\)，所求结论直接成立；\(K=0\) 几乎处处的情形也直接成立。其余情形对 \(s\) 积分，由 Tonelli 定理及分层表示得到
\[
 \begin{aligned}
 \int K_t(y)|f(x-y)|\,\mathrm dy
 &=\int_0^\infty\int_{E_s}|f(x-y)|\,\mathrm dy\,\mathrm ds\\
 &\leq Mf(x)\int_0^\infty m_d(E_s)\,\mathrm ds
 =\|K\|_1Mf(x).
 \end{aligned}
\]
积分三角不等式给出最后的断言。
\end{proof}

例如 \(K=\mathbf1_{B(0,1)}/\omega_d\) 给出原来的中心球平均。另取 \(K(x)=c\cdot(1+|x|)^{-d-1}\)，其中 \(c\) 使积分为 \(1\)，便得到支集无界的正核。归一化常数存在：单位球上函数有界，而在 \(2^j\leq|x|<2^{j+1}\) 上的积分至多为一个与 \(j\) 无关的常数乘以 \(2^{-j}\)。两种核都受同一个极大函数控制，尽管前者有跳跃，后者从所有远处位置取值。

\subsection{后续调和分析中的作用}
\label{b1:subsec:130404}

现在可以分清两种收敛论证的分工。平移连续性说明：对固定 \(f\)，积分核的集中足以使误差的 \(L^p\) 范数趋零。极大估计说明：如果在一个便于计算的稠密类上已知逐点收敛，那么所有尺度的误差可以同时控制，例外集的测度最终也趋零。前者是范数中的逼近，后者是逐点收敛的稳定性，不能把其中一个当作另一个的同义改写。

\ProseParagraphBreak

对径向非增核，上一命题的控制是逐点、全尺度的；需要直接使用可测极大函数时，可以先取可数尺度
\[
 K^*_{\mathbb Q}f(x)=\sup_{t\in\mathbb Q,\ t>0}|K_t*f(x)|.
\]
它是可测函数，并由 \(\|K\|_1Mf\) 控制。于是当 \(f\in L^1\)、\(\|K\|_1>0\) 时，
\[
 m_d\bigl(\{K^*_{\mathbb Q}f>\lambda\}\bigr)
 \leq\frac{5^d\|K\|_1}{\lambda}\|f\|_1;
\]
当 \(f\in L^p\)、\(1<p<\infty\) 时，
\[
 \|K^*_{\mathbb Q}f\|_p\leq\|K\|_1 C_{d,p}\|f\|_p.
\]
给定任意一列正数 \(t_n\downarrow0\)，同样可以对 \(\sup_n|K_{t_n}*f|\) 使用这些估计。若还要恢复所有实数尺度的极限，则须利用核族的结构或直接作逐点误差估计；下一章在 Lebesgue 点处采用后一条路线。

\ProseParagraphBreak

调和分析中还会遇到不能由绝对平均直接支配的算子。某些核靠正负抵消才有意义，先把核和函数全部取绝对值反而会破坏关键性质。届时覆盖选择、幅值分解与极大估计仍会使用，但要与抵消条件配合。程民德等的《实分析》\cite[第三章]{Cheng2008RealAnalysis}将覆盖、极大函数、插值及奇异积分放在相邻各节，提供了一条中文阅读路径；本卷先把覆盖与微分的基础做完，第二卷再进入这些算子的系统理论。
